% ID: SEL-D0082-OA-E01
% Tanda S082.
\textbf{a)} Despejamos las restricciones:
\[
 y\ge2-2x,\qquad y\ge\frac32x-\frac32,\qquad y\le\frac{x}{3}+2.
\]
El recinto es el triángulo cerrado situado por encima de las dos primeras rectas y por debajo de la tercera.
\begin{center}
\setlength{\unitlength}{1mm}
\begin{picture}(109,85)(-19,-22)
\put(-14,0){\vector(1,0){100}}\put(0,-17){\vector(0,1){73}}
\put(88,0){$x$}\put(1,56){$y$}
{\color{gray}
\qbezier(-8,50.4)(12,14.4)(32,-21.6)
\qbezier(-1,-28.35)(29.5,12.825)(60,54)
\qbezier(-10,33)(30,45)(70,57)}
\linethickness{0.45mm}
\qbezier(0,36)(10,18)(20,0)
\qbezier(20,0)(40,27)(60,54)
\qbezier(60,54)(30,45)(0,36)
\put(2,39){\scriptsize $(0,2)$}
\put(16,-7){\scriptsize $(1,0)$}
\put(61,54){\scriptsize $(3,3)$}
\end{picture}
\end{center}

\textbf{b)} Intersecando $y=2-2x$ con $y=x/3+2$ obtenemos $(0,2)$. Intersecando las dos rectas inferiores, $2-2x=3x/2-3/2$, resulta $(1,0)$. La intersección de $y=3x/2-3/2$ con $y=x/3+2$ es $(3,3)$. Cada punto satisface la restricción restante, de modo que son los tres vértices.

\textbf{c)} Al ser una función lineal sobre un polígono cerrado y acotado, basta evaluar los vértices:
\[
 F(0,2)=-2,\qquad F(1,0)=2,\qquad F(3,3)=3.
\]
El mínimo es $\boxed{-2}$ y se alcanza únicamente en $(0,2)$.
